\documentclass{article}
\usepackage{babel}
\selectlanguage{english}
\usepackage{amsmath}
\usepackage{amssymb}
\usepackage[colorlinks=true, linkcolor=blue]{hyperref}
\usepackage{xcolor}
\usepackage[top=2cm, bottom=2cm,
            left=1.5cm, right=1.5cm]{geometry}
%\usepackage{minted}
%\usemintedstyle{vs}
\usepackage{url}
\usepackage{graphicx}
\usepackage{cite}
\usepackage{ltablex}
\usepackage{fancyvrb}
\usepackage{hyperref}
\keepXColumns

\hypersetup{
    colorlinks=true,      % Turns off the boxes, colors the text instead
    linkcolor=blue,       % Color for internal links (sections, equations, figures)
    citecolor=teal,       % Color for citations (\cite{})
    urlcolor=cyan         % Color for web URLs (\url{})
}

% --- Inline style for code-level names (AthenaK variables, files, flags) ------
\definecolor{varbg}{HTML}{EEF1F7}
\definecolor{varfg}{HTML}{273C75}
\newcommand{\var}[1]{%
  {\setlength{\fboxsep}{1.5pt}%
   \colorbox{varbg}{\textcolor{varfg}{\texttt{\small #1}}}}}

% --- Run-in headers for the individual steps of a workflow -------------------
% Small caps + rule instead of \textbf, so they read as a level below
% \subsubsection rather than as another heading of the same weight.
\definecolor{stepfg}{HTML}{1A6B62}
\newcommand{\step}[1]{%
  \par\medskip\noindent
  {\color{stepfg}\rule[-0.15ex]{2pt}{1.5ex}\hspace{0.5em}\scshape #1}%
  \hspace{0.75em}\ignorespaces}

% Define short-commands for the variables
\newcommand{\gxx}{\var{adm\_gxx}}
\newcommand{\gxy}{\var{adm\_gxy}}
\newcommand{\gxz}{\var{adm\_gxz}}
\newcommand{\gyy}{\var{adm\_gyy}}
\newcommand{\gyz}{\var{adm\_gyz}}
\newcommand{\gzz}{\var{adm\_gzz}}
\newcommand{\vx}{\var{velx}}
\newcommand{\vy}{\var{vely}}
\newcommand{\vz}{\var{velz}}
\newcommand{\alp}{\var{z4c\_alpha}}
\newcommand{\betx}{\var{z4c\_betax}}
\newcommand{\bety}{\var{z4c\_betay}}
\newcommand{\betz}{\var{z4c\_betaz}}
\newcommand{\Bx}{\var{bcc1}}
\newcommand{\By}{\var{bcc2}}
\newcommand{\Bz}{\var{bcc3}}
\newcommand{\mr}{\mathrm}
\setlocalecaption{nil}{ref}{References}


\begin{document}
\section*{Driftcontrol in AthenaK}
\textit{This document collects the driftcontrol algorithms used in AthenaK to reduce
the black hole kick in various applications.}

\subsection*{Driftcontrol}
\step{Preliminaries}
In black-hole neutron star mergers the remnant black hole experiences a physical kick during merger which leads
to it moving away from the coordinate origin. For analysis purposes this is not ideal, since the matter distribution will follow the
compact object and in most cases the surface analysis for ejecta etc. is biased. In short simulations this might not make
a tremendous impact, but it certainly does in longer evolutions. As such we can seek to modify the gauge to pull the black hole
back to the desired target location and keep it there throughout the rest of the simulation. Due to the gauge-invariance of the system
this should not alter any physical properties. AthenaK implements four different methods that address this problem and which are
described below. Here, we first introduce the basic concept behind the driftcontrol.
\step{Basic concept} Consider that the drifting remnant black hole is followed by AthenaK's \var{CompactObjectTracker}, so we now it's
position and velocity over time. For a tracked black hole the velocity of the tracker is determined as $v_{\mathrm{tracker}}^i=-\beta^i$ and
the position is computed via $x_{\mathrm{tracker}}^i=x_{\mathrm{tracker,old}}^i+\Delta tv_{\mathrm{tracker}}^i$. AthenaK's driftcontrol never touches
the position, but instead adds a controlling term to the RHS of the shift evolution equation $\beta^i$. In particular, we seed something like
\begin{align}
  \partial_t\beta^i = ...-f_{\mathrm{corr}}g,
\end{align}
where $g$ is a Gaussian window, and $f_{\mathrm{corr}}=-u$ where $u$ is what each individual controller seeds based on the specified algorithm. For
a unit-mass particle we find
\begin{align}
  \Dot{x}^i&=-\beta^i,\\
  \Ddot{x}^i&=-\partial_t\beta^i=-u+f,
\end{align}
where $f$ is the unperturbed RHS of the shift evolution equation, i.e. the contribution from the $\Gamma$-driver, shift advection, and $\eta$-damping.
We therefore inject a control acceleration $u$ into the system that tries to control the drifting black hole. The algorithm is as follows:
\begin{enumerate}
  \item Measure the error of the current position to the target position $e=x-x_{\mathrm{target}}$.
  \item Compute the injected acceleration $u(e,v,\mathrm{state})$ defined by each control mechanism.
  \item Multiply that by a Gaussian window $g(r)$ which ideally is a certain factor larger than the black hole.
  \item Feed this into the evolution equation of the shift $\partial_t\beta+=u\cdot g$.
  \item The velocity will be automatically adjusted via $\Dot{x}^i=-\beta^i$.
  \item Steps 1-5 are repeated at every cycle at the end of every RK stage.
\end{enumerate}
Physically nothing changes, as the shift is pure gauge and does not enter the constraints. However, holding the black hole there introduces additional coordinate shear around
the black hole because the coordinates must keep stretching to keep the object there.
\step{Shared properties} The four driftcontrol algorithms that will be introduced below all share some common machinery inside of AthenaK:
\begin{itemize}
  \item \var{dc\_tracker\_index} selects which object shall be controlled (after merger both share the same object).
  \item \var{dc\_fixed\_x/y/z} is the target position we want the black hole at.
  \item \var{dc\_vel} is the velocity computed as the finite-difference of the tracker position between cycles, i.e. $v=\frac{x_n-x_{n-1}}{\Delta t}$, but it is not
  the tracker's own velocity. This velocity is clamped to $\pm$\var{dc\_vel\_cap}, which guards against bad interpolations or a tracker jump from injecting a
  huge derivative kick into the shift.
  \item The correction $f_{\mathrm{corr}}$ we feed into the shift is spatially constant and only the weight varies. With the Gaussian window we put the shear, where the
  peak of the Gaussian window is located (ideally at the black hole). \var{dc\_gaussian\_center} centers the Gaussian window on either \var{fixed} (the target position) or
  \var{tracker} (the tracked object). \var{dc\_damping\_scale} is the $\frac{1}{e}$ radius inside of $g(r)=\exp{\left(\frac{-r^2}{s^2}\right)}$, where $s$ is the damping scale.
\end{itemize}

\subsection*{Oscillator}
\step{Formulation} The first method that is implemented is the Oscillator method from THC (second-order). For this method, the injected acceleration into the shift evolution is given by
\begin{align}\label{Equ:oscillator}
  u=\frac{e}{\tau^2}+\frac{2\zeta v}{\tau},
\end{align}
which is a natural damped harmonic oscillator with frequency $\omega=\frac{1}{\tau}$, and where $e$ is the error of the current position with respect to the target.
It is unconditionally stable for $\zeta,\tau>0$. The parameters that can be tuned are
\begin{itemize}
  \item \var{dc\_damping\_time}: the restoring timescale $\tau=\frac{1}{\omega}$. A small $\tau$ means a stiff, and fast spring. This can lead to pullback speeds exceeding the
  speed-of-light. At critical damping the compact object settles in roughly 3-5$\tau$.
  \item \var{dc\_damping\_coeff}: the damping ratio $\zeta$. $\zeta<1$ represents an underdamped oscillator, $\zeta=1$ represents critical damping, and $\zeta>1$ overdamping.
  The default is critical damping.
\end{itemize}
\step{Caveat} Equ.~\eqref{Equ:oscillator} can be written as
\begin{align}
  \Ddot{x}+\frac{e}{\tau^2}+\frac{2\zeta}{\tau}\Dot{x}=f.
\end{align}
Setting $\Ddot{x}=v=0$ gives $e=f\tau^2$, which means with this method we retain a constant offset that scales with the square of the timescale. Halving the offset
costs a factor of $\sqrt{2}$ in stiffness. One can use it when the drift is mild and a constant offset is acceptable. It works better if it is turned off shortly after the
kick, such that the initial separation from the target is not too large yet.

\subsection*{PID}
\step{Formulation} The PID method is third-order and the injected acceleration is defined by
\begin{align}\label{Equ:pid}
  u=K_pe+K_i\int edt+K_dv,
\end{align}
where $u$ without the integral term and $K_p=\frac{1}{\tau^2}$, $K_d=\frac{2\zeta}{\tau}$ reproduce the harmonic oscillator of
Equ.~\eqref{Equ:oscillator} exactly. The integral term accumulates offset from the target and keeps pushing until the offset is gone,
which kills the problem the Oscillator method had. This extended method requires that the coefficients $K_i,K_p,K_d>0$ and $K_i<K_pK_d$. The
individual parameters that can be tuned are:
\begin{itemize}
  \item \var{dc\_Kp}: the spring constant $\omega_0=\sqrt{K_p}$, with the default being 1.0 such that $\tau=1$.
  \item \var{dc\_Kd}: the damping parameter $\zeta=\frac{K_d}{(2\sqrt{K_p})}$ (default is critically damped).
  \item \var{dc\_Ki}: the integral gain, i.e. how fast the accumulated offset is converted into a restoring force.
  \item \var{dc\_integral\_cap}: clamp on the integral $\int edt$.
\end{itemize}
\step{Caveat} The residual error in principle is now zero with this method. However, releasing the driftcontrol at one point is fatal, because the
restoring force used to keep the object at the desired location is then used to kick the black hole away again with the opposite force. The only option is to
keep the remnant at the location with PID on for the entire time.

\subsection*{Relaxation}
\step{Formulation} The Relaxation method is first-order and relaxes the local shift field toward a commanded value. It modifies the shift evolution equation in the form
\begin{align}
  \partial_t \beta^i-=\left(\frac{1}{\tau_r}\right)\left(\beta^i-\kappa e^i\right)g-G_i\Gamma^ig,
\end{align}
where $g$ again is the Gaussian window, $\tau_r$ is the timescale, and $G_i$ is a constant. If $\frac{1}{\tau_r}$ dominates the other gauge terms, then
$\beta^i\rightarrow\kappa e^i$ on timescale $\tau_r$. Then $\Dot{x}=-\beta=-\kappa e$ and the offset decays exponentially at rate $\kappa$. The tuning parameters are:
\begin{itemize}
  \item \var{dc\_kappa}: return rate $\kappa$ to the desired location.
  \item \var{dc\_relaxation\_time}: relaxation time $\tau_r$ similar to the Oscillator method.
  \item \var{dc\_gamma\_suppress}: fades out the $\Gamma$-driver's contribution to $\partial_t\beta^i$ inside the window (no suppression at $e=0$). In essence, the
  relaxation term and the $\Gamma$-driver are fighting and in principle, the further off target one is the more the driftcontrol should win.
\end{itemize}
\step{Caveat} At equilibrium $\beta=0$ we get a residual $e=\frac{f\tau_r}{\kappa}$, which is still biased but still much better than the Oscillator at the
same $\tau_r$. This is the only of the four methods that locally changes the shift gauge instead of nudging it. It is also not guarded against superluminal
coordinate speeds. However, it pins the object locally.

\subsection*{DOB (Disturbance observer)}
\step{Formulation} The DOB method is third-order and it uses the injected acceleration
\begin{align}
  \Hat{f}&=p+\omega_0v,\\
  u&=\omega_c^2e+2\zeta\omega_cv+\Hat{f},\\
  \Dot{p}&=\omega_0(\omega_c^2e+2\zeta\omega_cv).
\end{align}
We can re-arrange the equation for $u$ to $\omega_c^2e+2\zeta\omega_cv=u-\Hat{f}$, which leads to the identity $\Dot{p}=\omega_0(u-\Hat{f})$.
Differentiating $\Hat{f}$ leads to $\Dot{\Hat{f}}=\Dot{p}+\omega_0\Ddot{x}$. Substituting $\Dot{p}$ and the controller $\Ddot{x}=-u+f$ results in
$\Dot{\Hat{f}}=\omega_0(f-\Hat{f})$, i.e. the estimator does not depend on what the controller is doing. The first two terms in $u$ are essentially just what
the Oscillator method is doing. To understand the role of $\Hat{f}$ a bit better, we can split the injected acceleration in two parts
\begin{align}
  u&=u_{\mathrm{PD}}+\Hat{f},\ \text{PD=Proportional-Derivative}\\
  u_{\mathrm{PD}}&=\underbrace{\omega_c^2e}_{\text{proportional to e}}+\underbrace{2\zeta\omega_cv}_{\text{derivative of e}},
\end{align}
which if substituted into $\Ddot{x}=-u+f$ gives $\Ddot{x}=-u_{\mathrm{PD}}+(f-\Hat{f})$. $\Hat{f}$ therefore is trying to delete $f$ from the equation. Physically, the
observer is inferring the drive from the mismatch between the acceleration that wanted and the acceleration one actually gets. Put differently, we push with $-u$, watch the
puncture and then attribute the difference to the ambient gauge, i.e. $f$.
The tuning parameters are:
\begin{itemize}
  \item \var{dc\_omega\_c}: the analog to $\frac{1}{\tau}$ given in units $[\mathrm{M}^{-1}]$.
  \item \var{dc\_zeta}: the damping ratio, and identical to the Oscillator method. Critical damping, i.e. $\zeta=1.0$ means there is no overshoot.
  \item \var{dc\_omega\_o}: observer bandwidth, i.e. how fast $\Hat{f}$ tracks $f$. This improves the
  disturbance rejection without eating any stability margin. It is numerically limited by $\omega_0\Delta t<0.5$.
\end{itemize}
\step{Caveat} It is the preferred method over PID and the others, because it has no ceiling for instance on the integral as PID and if
the state is ever lost or wrong, it re-converges in $\approx\frac{1}{\omega_0}$.

\end{document}

